\documentclass[11pt]{article}
\usepackage[a4paper,margin=1in]{geometry}
\usepackage{amsmath,amssymb,mathtools}
\usepackage[colorlinks=true,linkcolor=blue,urlcolor=blue]{hyperref}
\usepackage[T1]{fontenc}
\usepackage[utf8]{inputenc}

\title{AKDE Current Math Formulas}
\author{}
\date{Generated from the active TypeScript implementation}

\newcommand{\formularefs}[1]{%
\par\smallskip
\noindent\footnotesize\textit{References: }#1
\par\smallskip\normalsize
}

\begin{document}
\maketitle

\section{Scope}

This document lists the formulas currently used by the AKDE implementation under
\texttt{kirigami/model/}. It reflects the active code paths, including clamps,
validation gates, and the currently selected minor-cut formula.

\section{Inputs and Symbols}

\begin{align*}
N &:= \text{edge count}, \qquad N \ge 3, \\
L &:= \text{base edge length (mm)}, \\
L_o &:= \text{outer polygon-face edge length on the perimeter (mm)}, \\
H &:= K_{\mathrm{tot}} = \text{vertical apex height (mm)}, \\
T &:= \text{material thickness (mm)}.
\end{align*}

The implementation validates
\[
H > 0,
\qquad
T > 0.
\]

\formularefs{Implementation authority: \texttt{kirigami/model/validation.ts}.}

\section{Derived Pyramid Geometry}

\subsection{Primary lengths and angles}

\begin{align}
R &= \frac{L}{2\sin\!\left(\frac{\pi}{N}\right)}, \\
s &= \sqrt{R^2 + H^2}, \\
\psi &= \operatorname{atan2}(H,R), \\
\kappa &= \frac{H}{R}.
\end{align}

\formularefs{For the regular-\(N\)-gon circumradius relation and non-right-triangle trigonometry, see OpenStax, ``\href{https://openstax.org/books/algebra-and-trigonometry-2e/pages/10-1-non-right-triangles-law-of-sines}{10.1 Non-right Triangles: Law of Sines}'' and ``\href{https://openstax.org/books/algebra-and-trigonometry-2e/pages/7-2-right-triangle-trigonometry}{7.2 Right Triangle Trigonometry}.''}

The apex face angle is
\begin{equation}
\eta =
\begin{cases}
2\arcsin\!\left(\min\!\left(1,\dfrac{R\sin(\pi/N)}{s}\right)\right), & s>0 \text{ and } R>0,\\[6pt]
0, & \text{otherwise.}
\end{cases}
\end{equation}

\formularefs{The \(\eta\) expression is the law-of-sines reduction for the lateral face isosceles triangle; see OpenStax, ``\href{https://openstax.org/books/algebra-and-trigonometry-2e/pages/10-1-non-right-triangles-law-of-sines}{10.1 Non-right Triangles: Law of Sines}.'' AKDE-specific specialization: \texttt{kirigami/model/geometry.ts}.}

The discrete angle defect and uniform molecule angle are
\begin{align}
\delta_{\mathrm{apex}} &= 2\pi - N\eta, \\
\theta &= \frac{\delta_{\mathrm{apex}}}{N}, \\
\tau &= \frac{2\pi}{N}.
\end{align}

\formularefs{For angle defect as \(2\pi\) minus the sum of incident angles, see Keenan Crane, \emph{Discrete Differential Geometry: An Applied Introduction}, §5.5--5.6, especially the angle-defect definition and discrete Gauss--Bonnet discussion: \href{https://www.cs.cmu.edu/~kmcrane/Projects/DDG/paper.pdf}{CMU notes PDF}. For the kirigami reduction \(\theta=\delta_{\mathrm{apex}}/N\) and closure step \(\tau=2\pi/N\), see Liu et al., ``\href{https://static1.squarespace.com/static/5783b6f903596e5098f3fce8/t/5dd2e95aee053940522eab1a/1574103423054/DETC2019-97557.pdf}{Programmable Kirigami},'' Eq.~(1) and Figure~2.}

The molecule width is
\begin{equation}
w = 2s\sin\!\left(\frac{\theta}{2}\right).
\end{equation}

\formularefs{The paper defines \(W(i,j)\) as the edge-molecule width and Figure~2 gives the symmetric trapezoid geometry; the chord specialization used here is an AKDE uniform-fan derivation built on Liu et al., ``\href{https://static1.squarespace.com/static/5783b6f903596e5098f3fce8/t/5dd2e95aee053940522eab1a/1574103423054/DETC2019-97557.pdf}{Programmable Kirigami},'' Figure~2 and Eq.~(2). Implementation authority: \texttt{kirigami/model/geometry.ts}.}

\subsection{Major cut and molecule measures}

The physical major-cut radius is
\[
r_{\mathrm{phys}} = \frac{T}{\sin(\theta/2)}.
\]

The implementation applies a visualization clamp:
\begin{equation}
r_{\mathrm{apex}} =
\begin{cases}
\min\!\left(\dfrac{T}{\sin(\theta/2)},\,0.4\,s\right), & T>0 \text{ and } \theta>0,\\[8pt]
0, & \text{otherwise.}
\end{cases}
\end{equation}

The molecule end-leg length is
\begin{equation}
D = 2s\tan\!\left(\frac{\theta}{2}\right)
= \frac{w}{\cos(\theta/2)},
\qquad
0 < \theta < \pi.
\end{equation}

\formularefs{Major/minor cuts and molecule geometry are introduced in Liu et al., ``\href{https://static1.squarespace.com/static/5783b6f903596e5098f3fce8/t/5dd2e95aee053940522eab1a/1574103423054/DETC2019-97557.pdf}{Programmable Kirigami},'' Figure~2 and the paragraph introducing major and minor cuts. The explicit formulas for \(r_{\mathrm{apex}}\), the \(0.4s\) clamp, and \(D\) are AKDE implementation formulas in \texttt{kirigami/model/geometry.ts}.}

\subsection{Dihedral angle}

Let
\[
\phi = -\frac{\pi}{2},
\qquad
\Delta = \frac{2\pi}{N},
\qquad
\mathbf{a} = (0,0,H),
\]
and define base vertices
\[
\mathbf{b}(\alpha) = \bigl(R\cos\alpha,\; R\sin\alpha,\; 0\bigr).
\]

Then
\begin{align*}
\mathbf{b}_i &= \mathbf{b}(\phi), \\
\mathbf{b}_{\mathrm{prev}} &= \mathbf{b}(\phi-\Delta), \\
\mathbf{b}_{\mathrm{next}} &= \mathbf{b}(\phi+\Delta).
\end{align*}

The adjacent face normals used by the code are
\begin{align}
\mathbf{n}_1 &= (\mathbf{b}_{\mathrm{prev}}-\mathbf{b}_i) \times (\mathbf{a}-\mathbf{b}_i), \\
\mathbf{n}_2 &= (\mathbf{a}-\mathbf{b}_i) \times (\mathbf{b}_{\mathrm{next}}-\mathbf{b}_i).
\end{align}

If either normal has near-zero length, the implementation returns $\gamma=\pi$.
Otherwise,
\begin{align}
\alpha_{\mathrm{norm}} &=
\arccos\!\left(
\operatorname{clamp}_{[-1,1]}
\frac{\mathbf{n}_1\cdot\mathbf{n}_2}{\lVert \mathbf{n}_1\rVert \lVert \mathbf{n}_2\rVert}
\right), \\
\gamma &= \pi - \alpha_{\mathrm{norm}}.
\end{align}

\formularefs{For cross products and normals to the plane spanned by two vectors, see OpenStax, ``\href{https://openstax.org/books/calculus-volume-3/pages/2-4-the-cross-product}{2.4 The Cross Product}.'' The specific adjacent-face construction and \(\gamma=\pi-\alpha_{\mathrm{norm}}\) convention are AKDE implementation choices in \texttt{kirigami/model/geometry.ts}.}

\section{Minor-Cut Formulas Currently Active}

The active implementation selects \texttt{MINOR\_CUT\_FORMULA = "fold-reach"}.

First, the fold-clearance depth used by constraint $C6$ is
\begin{equation}
d_{\mathrm{clear}} =
\begin{cases}
T\tan(\gamma/2), & T>0,\; 0<\gamma/2<\pi/2,\\[4pt]
0, & \text{otherwise.}
\end{cases}
\end{equation}

The active minor-cut length is
\begin{equation}
\ell = \sqrt{\left(\frac{w}{2}\right)^2 + r_{\mathrm{apex}}^2}.
\end{equation}

This is the same as
\[
\ell = \operatorname{hypot}\!\left(\frac{w}{2},\,r_{\mathrm{apex}}\right).
\]

\formularefs{Liu et al. state that the minor cut depends on ``the angle between adjacent mesh faces and the width of the hinge'' and refer to Figure~2, but they do not give a closed-form equation; see ``\href{https://static1.squarespace.com/static/5783b6f903596e5098f3fce8/t/5dd2e95aee053940522eab1a/1574103423054/DETC2019-97557.pdf}{Programmable Kirigami},'' the paragraph immediately before Figure~2. The formulas \(d_{\mathrm{clear}}=T\tan(\gamma/2)\) and \(\ell=\sqrt{(w/2)^2+r_{\mathrm{apex}}^2}\) are AKDE-specific constructions in \texttt{kirigami/model/geometry.ts}; the current branch is \texttt{MINOR\_CUT\_FORMULA = "fold-reach"}.}

\section{Constraint Equations}

The code uses a tolerance
\[
\varepsilon = 10^{-10}.
\]

\subsection{$C1$: angle closure}

\begin{align}
\rho_{C1} &= \left|N\theta + N\eta - 2\pi\right|, \\
C1 &\text{ is satisfied iff } \rho_{C1} < \varepsilon.
\end{align}

\formularefs{Directly from Liu et al., Eq.~(1): ``\href{https://static1.squarespace.com/static/5783b6f903596e5098f3fce8/t/5dd2e95aee053940522eab1a/1574103423054/DETC2019-97557.pdf}{Programmable Kirigami}.''}

\subsection{$C2$: vector closure}

The phase sequence is
\[
\Phi_n = n\tau,
\qquad n=0,1,\dots,N-1.
\]

The closure sum is
\begin{align}
S_x &= \sum_{n=0}^{N-1} w\cos(\Phi_n), \\
S_y &= \sum_{n=0}^{N-1} w\sin(\Phi_n),
\end{align}
with residual
\begin{align}
\rho_{C2} &= \sqrt{S_x^2 + S_y^2}, \\
C2 &\text{ is satisfied iff } \rho_{C2} < \varepsilon.
\end{align}

\formularefs{Directly from Liu et al., Eq.~(2) and the closed-polygon condition explained immediately below it in ``\href{https://static1.squarespace.com/static/5783b6f903596e5098f3fce8/t/5dd2e95aee053940522eab1a/1574103423054/DETC2019-97557.pdf}{Programmable Kirigami}.''}

\subsection{$C3$: molecule angle range}

\begin{align}
\rho_{C3} &= \max\!\left(0,\; |\theta| - (\pi-\varepsilon)\right), \\
C3 &\text{ is satisfied iff } |\theta| < \pi-\varepsilon.
\end{align}

\formularefs{Directly from Liu et al., inequality~(3) in ``\href{https://static1.squarespace.com/static/5783b6f903596e5098f3fce8/t/5dd2e95aee053940522eab1a/1574103423054/DETC2019-97557.pdf}{Programmable Kirigami}.''}

\subsection{$C4$: non-negative width}

\begin{align}
\rho_{C4} &=
\begin{cases}
-w, & w<0,\\
0, & w\ge 0,
\end{cases} \\
C4 &\text{ is satisfied iff } w \ge -\varepsilon.
\end{align}

\formularefs{Directly from Liu et al., inequality~(4) in ``\href{https://static1.squarespace.com/static/5783b6f903596e5098f3fce8/t/5dd2e95aee053940522eab1a/1574103423054/DETC2019-97557.pdf}{Programmable Kirigami}.''}

\subsection{$C5$: fold overlap}

\begin{align}
\rho_{C5} &= \max(0,\,w-L), \\
C5 &\text{ is satisfied iff } \rho_{C5} \le \varepsilon,
\end{align}
equivalently
\[
w \le L.
\]

\formularefs{This is an AKDE validity proxy, not a closed-form condition from the paper. The paper states that the molecule must be valid without overlapping; AKDE instantiates that idea as \(w \le L\). Implementation authority: \texttt{kirigami/model/constraints.ts}; conceptual context from Liu et al., the text introducing inequalities~(3)--(4).}

\subsection{$C6$: cut vs.\ dihedral}

The code computes
\begin{align}
d_{\mathrm{relief}} &=
\begin{cases}
T\tan(\gamma/2), & 0<\gamma/2<\pi/2-10^{-12},\\[4pt]
\infty, & \text{otherwise,}
\end{cases} \\
d_{\mathrm{avail}} &= s-r_{\mathrm{apex}}, \\
\rho_{C6} &= \max(0,\,d_{\mathrm{relief}} - d_{\mathrm{avail}}).
\end{align}

Constraint $C6$ is satisfied iff
\[
\rho_{C6} \le \varepsilon,
\]
which corresponds to the intended inequality
\[
T\tan(\gamma/2) \le s-r_{\mathrm{apex}}.
\]

\formularefs{This is an AKDE validity proxy derived from the minor-cut clearance concept, not a closed-form inequality from Liu et al. Conceptual source: Liu et al., major/minor-cut discussion around Figure~2. Implementation authority: \texttt{kirigami/model/constraints.ts}.}

\section{Rendered Net Geometry}

The exported SVG pattern uses the same derived geometry plus the optional perimeter
input $L_o$.

\subsection{Outer polygon span}

The implementation first chooses
\[
L_o^\ast =
\begin{cases}
L_o, & L_o>0,\\
L, & \text{otherwise.}
\end{cases}
\]

Then it computes
\begin{align}
\alpha_{\max} &= \frac{\tau}{2} - 10^{-3}, \\
\alpha_o &= \min\!\left(
\max\!\left(
\arcsin\!\left(\min\!\left(1,\frac{L_o^\ast}{2s}\right)\right),
10^{-4}
\right),
\alpha_{\max}
\right).
\end{align}

\formularefs{This is an AKDE rendering construction for the outer perimeter chord span. The triangle relation behind \(\arcsin(L_o^\ast/(2s))\) comes from the same law-of-sines / chord geometry used above; see OpenStax, ``\href{https://openstax.org/books/algebra-and-trigonometry-2e/pages/10-1-non-right-triangles-law-of-sines}{10.1 Non-right Triangles: Law of Sines}.'' Implementation authority: \texttt{kirigami/model/pattern.ts}.}

Here $\alpha_o$ is the half-span of each polygon outer base on the radius-$s$ circle.

\subsection{Rendered molecule outer chord}

Because the perimeter alternates polygon and molecule spans, the rendered molecule
chord becomes
\begin{equation}
w_{\mathrm{rendered}} = 2s\sin\!\left(\max\!\left(0,\frac{\tau}{2}-\alpha_o\right)\right).
\end{equation}

\formularefs{AKDE rendering-specific chord formula from \texttt{kirigami/model/pattern.ts}; it follows from circle-chord geometry after allocating perimeter span \(\alpha_o\) to the polygon outer base.}

\subsection{Minor cut passed to rendering}

The renderer uses
\begin{equation}
\ell_{\mathrm{render}} =
\operatorname{computeMinorCutLength}\!\left(\gamma,\,w_{\mathrm{rendered}},\,T,\,\theta,\,r_{\mathrm{apex}}\right).
\end{equation}

With the active \texttt{fold-reach} branch, this reduces to
\[
\ell_{\mathrm{render}} =
\sqrt{\left(\frac{w_{\mathrm{rendered}}}{2}\right)^2 + r_{\mathrm{apex}}^2}.
\]

\formularefs{AKDE rendering specialization of the active \texttt{fold-reach} branch from \texttt{kirigami/model/geometry.ts} and \texttt{kirigami/model/pattern.ts}.}

\section{Default Input Case}

The default state is
\[
N=6,
\qquad
L=100~\mathrm{mm},
\qquad
L_o=100~\mathrm{mm},
\qquad
T=1~\mathrm{mm},
\]
with
\[
R = \frac{L}{2\sin(\pi/N)},
\qquad
H = R.
\]

\end{document}
